% Unite: papier 24, introduction et section 1; traduction francaise depuis la source allemande auditee.
\setcounter{footnote}{0}
\providecommand{\frakS}{\mathfrak S}
\providecommand{\fraka}{\mathfrak a}
\providecommand{\frakb}{\mathfrak b}
\providecommand{\frakc}{\mathfrak c}
\providecommand{\frakg}{\mathfrak g}
\providecommand{\frakG}{\mathfrak G}
\providecommand{\frakh}{\mathfrak h}
\providecommand{\frakm}{\mathfrak m}
\providecommand{\frakM}{\mathfrak M}
\providecommand{\frakn}{\mathfrak n}
\providecommand{\frako}{\mathfrak o}
\providecommand{\frakp}{\mathfrak p}
\providecommand{\frakq}{\mathfrak q}
\providecommand{\frakr}{\mathfrak r}
\providecommand{\barfrakm}{\overline{\mathfrak m}}

\section*{24. Théorie de l'élimination et théorie générale des idéaux}
\begin{center}
Math. Ann. 90 (1923), p. 229--261\\[0.9em]
{\Large\bfseries Théorie de l'élimination et théorie générale des idéaux.}\\[1.0em]
Par\\[0.45em]
Emmy Noether à Göttingen.\\[0.6em]
\rule{4em}{0.4pt}
\end{center}

Il s'agit, dans ce qui suit, d'inscrire la théorie de l'élimination -- dans la forme arithmétique que j'ai donnée à l'exposé de Hentzelt\footnote{Hentzelt, Kurt: \emph{Zur Theorie der Polynomideale und Resultanten}. Édité par E. Noether. Math. Ann. 88 (1922), p. 53--79; cité H.-N.} et qui peut se désigner comme une théorie des idéaux dans l'anneau de polynômes -- dans la théorie générale des idéaux,\footnote{Noether, E.: \emph{Idealtheorie in Ringbereichen}, Math. Ann. 83 (1921), p. 24--66; cité \emph{Idealtheorie}.} et en même temps de donner une refondation des parties relatives aux zéros. Les notions fondamentales des deux théories sont brièvement réunies au \S{} 1, de même que celles de la théorie des corps,\footnote{Steinitz, E.: \emph{Algebraische Theorie der Körper}, J. f. M. 137 (1910), p. 167--309; cité Steinitz.} de sorte que je pourrai m'y référer ici.

L'insertion et la refondation se groupent autour de quatre questions: formulation arithmétique de la notion de dimension; théorie des zéros; rapport des théorèmes de décomposition pour les normes et les diviseurs élémentaires avec ceux de la théorie générale des idéaux; caractérisation des idéaux premiers et primaires par la norme et la forme des diviseurs élémentaires.

La formulation arithmétique de la \emph{notion de dimension} (\S{} 4) se fait par des chaînes d'idéaux premiers, équivalent arithmétique du fait que, sur les surfaces, il y a des courbes, et sur les courbes, des points. Cette formulation arithmétique, dont on démontre l'identité avec la définition par paramètres de la théorie de l'élimination, se transpose avec une légère modification à des anneaux arbitraires et y donne, conjointement avec la transposition de certaines parties des autres questions, une vue exacte de la structure des idéaux, comme je le montrerai ailleurs.

La question des \emph{zéros} est abordée chez H.-N., comme d'ordinaire, directement pour l'idéal donné, à savoir par retour à l'élimination successive, ce qui n'évite pas entièrement le caractère de vérification. À l'opposé se trouve la voie toujours suivie pour les polynômes d'une seule variable: on représente le polynôme donné comme produit de puissances de polynômes irréductibles (polynômes primaires), et l'on construit, pour chaque polynôme irréductible pris isolément, le corps des zéros comme un corps isomorphe au corps résiduel. Le degré du polynôme irréductible donne le degré de ce corps; mais le degré du polynôme primaire, dont les zéros coïncident avec ceux du polynôme irréductible, donne la dimension de l'anneau quotient par ce polynôme primaire. Si l'on passe, pour chaque polynôme irréductible, au corps de Galois, on obtient la décomposition en facteurs linéaires; et ainsi, pour le polynôme primitivement donné, la question des zéros, de la décomposition en facteurs linéaires et de la multiplicité est réglée. De manière tout à fait correspondante, on passe ici (\S{} 3) de l'anneau quotient par un idéal premier à son corps résiduel, et l'on construit un corps des zéros qui lui est isomorphe. Ce corps des zéros contient un sous-corps engendré par adjonction d'indéterminées -- par exemple $x_{i+1},\ldots,x_n$; le degré de la norme donne le degré par rapport à ce sous-corps; en même temps, le passage au corps de Galois fournit la décomposition de la forme des diviseurs élémentaires et par conséquent de la norme en facteurs linéaires. Pour l'idéal primaire associé, les zéros sont cependant les mêmes; la décomposition est également donnée avec lui, puisque la norme devient une puissance de la forme des diviseurs élémentaires de l'idéal premier (\S{} 2); le degré de la norme donne la dimension de l'anneau quotient par l'idéal primaire relativement au sous-corps issu des indéterminées.

Ainsi la question des zéros d'un idéal arbitraire est ramenée au rapport des théorèmes de décomposition pour les normes et les diviseurs élémentaires avec ceux de la théorie générale des idéaux (\S{} 5). Au fait que les zéros de l'idéal se composent à partir de ceux de ses divers idéaux premiers associés correspond la décomposition de la norme en plus grands facteurs primaires, qui deviennent égaux à des puissances de la forme des diviseurs élémentaires des divers idéaux premiers associés;\footnote{Ce parallélisme de la théorie de l'élimination avec la théorie générale des idéaux n'est pas réalisé dans la théorie de l'élimination de Kronecker; voir à ce sujet H.-N., note *).} ainsi la décomposition en facteurs linéaires est accomplie pour la norme en général. Mais la multiplicité reçoit chez H.-N. son interprétation au moyen des idéaux fondamentaux et de leurs composantes; le degré d'un plus grand facteur primaire de la norme devient égal à la dimension -- après adjonction de $x_{i+1},\ldots,x_n$ -- du quotient de l'idéal fondamental de $(i-1)$-ième étage par une composante de l'idéal fondamental de $i$-ième étage. Or l'idéal fondamental de $(i-1)$-ième étage est démontré identique à la composante isolée de la décomposition qui est définie de manière unique par tous les idéaux premiers de dimension supérieure à la $(n-i)$-ième; et une composante de l'idéal fondamental de $i$-ième étage est démontrée identique à la composante isolée de la décomposition qui est définie par l'idéal premier correspondant au facteur de la norme et par tous les idéaux premiers de dimension supérieure. Et le degré du facteur correspondant de la norme devient -- après adjonction de $x_{i+1},\ldots,x_n$ -- égal à la dimension du sous-anneau de l'anneau quotient par cette composante formé des classes modulo celle-ci qui sont représentées par des éléments de l'idéal fondamental de $(i-1)$-ième étage.

La caractérisation des idéaux premiers et des idéaux primaires propres (\S{} 6) résulte comme conséquence directe de la considération des corps résiduels correspondants. Si le domaine des coefficients initial est un corps parfait, alors les idéaux premiers ont pour norme et forme des diviseurs élémentaires des polynômes irréductibles, les idéaux primaires propres des polynômes primaires propres, et réciproquement. Si le domaine des coefficients est un corps imparfait, on ne peut donner que des conditions soit nécessaires, soit suffisantes; la norme d'un idéal premier peut devenir un polynôme primaire propre, et la forme des diviseurs élémentaires d'un idéal primaire propre peut devenir un polynôme irréductible, comme le montrent des exemples. Plus précisément, en caractéristique $p$, la norme d'un idéal premier devient la $p^g$-ième puissance ($g\geqq0$) de sa forme des diviseurs élémentaires; tandis que, pour des idéaux primaires propres dont la forme des diviseurs élémentaires est un polynôme irréductible, la norme peut devenir une puissance quelconque, comme le montrent encore des exemples. On considère enfin les idéaux premiers absolus (\S{} 7), c'est-à-dire ceux qui restent des idéaux premiers lorsque le domaine des coefficients est étendu en un corps algébriquement clos. Pour des idéaux premiers absolus à coefficients dans un corps de nombres algébriques (fini), la caractérisation conduit à la transposition d'un théorème d'abord énoncé par A. Ostrowski pour les polynômes absolument irréductibles: la propriété d'être un idéal premier est conservée modulo chaque idéal premier du corps de nombres, sauf tout au plus un nombre fini d'exceptions.

Signalons encore l'analogie de cette dernière question avec la théorie des idéaux dans les corps de nombres algébriques. Comme forme des diviseurs élémentaires d'un tel idéal, il faut prendre le plus petit entier rationnel divisible par l'idéal (note 10), qui devient donc toujours un nombre premier pour un idéal premier, tandis que, réciproquement, un idéal devient premier dès que sa norme est un nombre premier; les preuves aussi se déroulent tout à fait parallèlement. Les exemples mentionnés ci-dessus montrent donc que, dans le cas d'un corps imparfait, apparaît l'analogue des idéaux premiers de degré supérieur et des idéaux de ramification, tandis que, dans le cas d'un corps parfait, il n'y a que des idéaux premiers du premier degré et aucun idéal de ramification.

\subsection*{\S{} 1. Notions fondamentales de la théorie de l'élimination, de la théorie générale des idéaux et de la théorie des corps}

\textbf{1. L'anneau.} Soit $\bar{\frakS}$ le domaine d'intégrité des polynômes en $y_1,\ldots,y_n$ à coefficients dans un corps $P$ défini abstraitement. Que les $y$ soient soumis à une transformation ayant pour coefficients des indéterminées $u_{\mu\nu}$:\footnote{La transformation est, en vue du \S{} 3, etc., un peu plus générale que chez H.-N.; tous les résultats qui s'y trouvent sont par là conservés à plus forte raison.}
\begin{align}
 y_1&=u_{11}x_1+\cdots+u_{1n}x_n,\quad \ldots,\quad
 y_n=u_{n1}x_1+\cdots+u_{nn}x_n, \tag{1}
\end{align}
avec l'inverse
\begin{align}
 x_1&=t_{11}y_1+\cdots+t_{n1}y_n,\quad \ldots,\quad
 x_n=t_{1n}y_1+\cdots+t_{nn}y_n, \tag{1'}
\end{align}
et que les $u_{\mu\nu}$, ou -- ce qui est équivalent en raison de leur exprimabilité rationnelle mutuelle -- les $t_{\mu\nu}$, soient adjoints à $P$, de sorte que naisse le domaine des coefficients $P(u)=P(t)$; $\frakS$ désignera le domaine d'intégrité des polynômes en $x_1,\ldots,x_n$ à coefficients dans $P(u)$. Les domaines $\frakS$ et $\bar{\frakS}$ sont à la base de la théorie de l'élimination; on considère des idéaux $\fraka,\frakb,\ldots$ dans $\frakS$ et $\bar{\fraka},\bar{\frakb},\ldots$ dans $\bar{\frakS}$. Un idéal dans un anneau quelconque est alors défini, comme d'habitude, par le fait qu'il contient, avec $\alpha$ et $\beta$, aussi $\alpha-\beta$, et, avec $\alpha$, aussi $\lambda\alpha$, où $\lambda$ désigne un élément arbitraire de l'anneau; $\frako$ désignera toujours l'idéal unité, composé de tous les éléments.

\textbf{2. Idéaux transformés.} Un idéal $\frakm$ dans $\frakS$ s'appelle transformé s'il naît d'un idéal $\barfrakm$ dans $\bar{\frakS}$ au moyen de (1), après adjonction des $u_{\mu\nu}$, ou des $t_{\mu\nu}$. Ainsi, lorsque $\bar f(y)$ ou $\bar g(y)$ parcourt tous les polynômes de $\barfrakm$, $\frakm$ se compose de toutes les combinaisons linéaires
\[
 f(x)=\sum U_i(u)\bar f_i(y)=\sum U_i(u)f_i(x)
\]
ou encore
\[
 g(x)=\sum T_i(t)\bar g_i(y)=\sum T_i(t)g_i(x);
\]
en particulier, tous les $f_i(x)=\bar f_i(y)$ y sont contenus, ou, ce qui revient au même, tous les $g_i(x)=\bar g_i(y)$. Comme $f(x)$, respectivement $g(x)$, ne sont déterminés qu'à des grandeurs de $P(u)=P(t)$ près, les $U_i,T_i$ peuvent être supposés, sans restriction de généralité, être des produits de puissances des $u$, respectivement des $t$. Les $f(x),g(x)$ passent de nouveau en des polynômes de $\frakm$ lorsque $U,T$ sont remplacés par des produits de puissances quelconques, donc en particulier par permutation des colonnes de (1), respectivement des lignes de (1'); ainsi $\frakm$ contient, avec n'importe quel polynôme, aussi tous ceux qui s'en déduisent par permutation des $x$, pourvu seulement que les permutations correspondantes soient effectuées dans les coefficients dépendant des $u$, respectivement des $t$. Tous les idéaux qui apparaîtront dans la suite seront -- sauf mention expresse du contraire -- supposés transformés. Les idéaux non transformés passent, par composition de (1) avec
\[
 x_i=v_{i1}z_1+\cdots+v_{in}z_n
\]
en idéaux transformés dans l'anneau de polynômes en $z$ à coefficients dans $P(u,v)$, tandis que, pour les idéaux transformés, cette composition revient seulement à remplacer les $u_{\mu\nu}$ par des combinaisons bilinéaires des $u,v$.

\textbf{3. Notation.} Par $f^{(i)},g^{(i)},\ldots,a^{(i)},b^{(i)},\ldots$ on entendra constamment des polynômes de $\frakS$ qui sont libres de $x_1,\ldots,x_{i-1}$.

\textbf{4. Idéaux fondamentaux.} À chaque idéal $\frakm$ sont associés $n$ idéaux fondamentaux $\frakg_0,\frakg_1,\ldots,\frakg_{n-1}$ par la convention suivante: l'idéal fondamental de $(i-1)$-ième étage $\frakg_{i-1}$ contient tous et seulement les polynômes $G(x)$ pour lesquels il existe un polynôme $b^{(i)}\ne 0$ -- en général variable avec $G(x)$ -- tel que $b^{(i)}G(x)\equiv 0\pmod{\frakm}$. L'existence de la base d'idéal garantit également celle d'au moins un $B^{(i)}$, fixe pour tous les $G(x)$, de sorte que
\begin{equation}
 B^{(i)}\frakg_{i-1}\equiv0\pmod{\frakm},\qquad B^{(i)}\ne0 . \tag{2}
\end{equation}
Les idéaux fondamentaux d'idéaux transformés deviennent eux-mêmes des idéaux transformés (H.-N., théorème VI). L'idéal fondamental de $i$-ième étage $\frakg_i$ est aussi défini comme l'idéal dans lequel $\frakm$ passe lorsqu'on adjoint $x_{i+1},\ldots,x_n$ à $P(u)$, si l'on se restreint -- ce qui alors ne signifie aucune restriction de généralité -- aux polynômes entiers en $x_{i+1},\ldots,x_n$; $\frakg_0$ devient toujours égal à l'idéal unité $\frako$.

\textbf{5. Représentation modulaire des idéaux, diviseurs élémentaires et normes individuelles.} Si l'on conçoit chaque polynôme $f(x)$ comme une forme linéaire dans les produits de puissances $\xi$ de $x_1,\ldots,x_{i-1}$, avec des polynômes $a^{(i)}$ comme coefficients, alors l'idéal $\frakm$ passe en un module $\frakM_{i-1}$ de formes linéaires relativement au domaine des $a^{(i)}$; c'est-à-dire que $\frakM_{i-1}$ contient, avec deux formes linéaires quelconques, aussi leur différence, et, avec une forme linéaire $l(\xi)$, aussi $a^{(i)}l(\xi)$ pour tout $a^{(i)}$. On appelle module fondamental d'un tel module $\frakM$ la totalité des formes linéaires $g(\xi)$ pour lesquelles $b^{(i)}g(\xi)\equiv0\pmod{\frakM}$, $b^{(i)}\ne0$; l'idéal fondamental $\frakg_{i-1}$ passe donc, dans la représentation modulaire, au module fondamental $\frakG_{i-1}$ de $\frakM_{i-1}$. $\frakM_{i-1}$ et $\frakG_{i-1}$ dépendent d'une infinité d'indéterminées $\xi$, de telle manière que, dans chaque forme linéaire prise isolément, seulement un nombre fini de ces indéterminées interviennent. $\frakM_{i-1}$ possède la propriété caractéristique de n'avoir, après adjonction de $x_{i+1},\ldots,x_n$ à $P(u)$, qu'un nombre fini de diviseurs élémentaires différents de l'unité; c'est-à-dire que $\frakG_{i-1}$ et $\frakM_{i-1}$ admettent, après cette adjonction, une représentation par base -- où $\zeta$ désigne de nouvelles indéterminées liées aux $\xi$ par une transformation linéaire réversible, entière en $x_i$, de telle manière que dans $\zeta_i=r_i(\xi)$ n'interviennent chaque fois qu'un nombre fini de ces indéterminées --
\begin{align}
\frakG_{i-1}&=(\zeta_0,\zeta_1,\ldots,\zeta_\sigma,\zeta_{\sigma+1},\ldots,\zeta_{\sigma+\nu},\ldots), \notag\\
\frakM_{i-1}&=(E^{(i)}\zeta_0,E_1^{(i)}\zeta_1,\ldots,E_\sigma^{(i)}\zeta_\sigma,
\zeta_{\sigma+1},\ldots,\zeta_{\sigma+\nu},\ldots), \tag{3}
\end{align}
où chaque $E_\mu^{(i)}$ divise le précédent (H.-N. (24) et (28)).

Le plus haut diviseur élémentaire $E^{(i)}$ est aussi défini comme le plus grand diviseur commun -- au sens polynomial -- de tous les $B^{(i)}$ qui apparaissent dans (2); il peut être supposé entier et primitif en $x_{i+1},\ldots,x_n$, et devient alors régulier en $x_i$, c'est-à-dire que $E^{(i)}$ contient le terme $x_i^\lambda$ avec coefficient non nul, s'il est de degré $\lambda$ en les $x$.

Le produit $R^{(i)}$ de tous les diviseurs élémentaires qui apparaissent dans (3) est appelé norme de $\frakG_{i-1}$ par rapport à $\frakM_{i-1}$, ou encore norme individuelle de l'idéal $\frakm$; en symboles, et en tenant compte du fait que $\frakm$ passe en $\frakg_i$ lors de l'adjonction de $x_{i+1},\ldots,x_n$:
\begin{equation}
 R^{(i)}=E^{(i)}E_1^{(i)}\cdots E_\sigma^{(i)}
 =N(\frakG_{i-1}\mid\frakM_{i-1})=N(\frakg_{i-1}\mid\frakg_i).
\end{equation}
Comme le montre (3), $R^{(i)}$ devient, après adjonction de $x_{i+1},\ldots,x_n$, égal au déterminant de la substitution de passage de $\frakG_{i-1}$ à $\frakM_{i-1}$; et le degré de $R^{(i)}$ donne la dimension, relativement à $P(u,x_{i+1},\ldots,x_n)$, du quotient $\frakG_{i-1}/\frakM_{i-1}$, donc de $\frakg_{i-1}/\frakg_i$; $R^{(i)}$ peut être supposé entier et primitif en $x_{i+1},\ldots,x_n$, et devient alors régulier en $x_i$. $R^{(i)}$ se calcule comme le plus grand diviseur commun -- au sens polynomial -- de tous les déterminants à $\varrho$ lignes d'un module $\frakM^*_{i-1}$, toujours existant, de rang $\varrho$, ayant la propriété que le système quotient $\frakG^*_{i-1}/\frakM^*_{i-1}$ est isomorphe à $\frakG_{i-1}/\frakM_{i-1}$, $\frakG^*_{i-1}$ étant entendu comme le module fondamental de $\frakM^*_{i-1}$ (H.-N., théorème VII et \S{} 5).

\textbf{6. Forme des diviseurs élémentaires et norme (forme résultante) des idéaux.} Le produit $E_\frakm=E^{(1)}E^{(2)}\cdots E^{(n)}$ de tous les plus hauts diviseurs élémentaires est appelé forme des diviseurs élémentaires, et le produit $R_\frakm=R^{(1)}R^{(2)}\cdots R^{(n)}$ de toutes les normes individuelles, norme (forme résultante)\footnote{Chez H.-N. seule la désignation forme résultante est employée; la désignation norme correspond aux propriétés arithmétiques.} de $\frakm$. La forme des diviseurs élémentaires et la norme sont divisibles par l'idéal; la norme est divisible par la forme des diviseurs élémentaires, et une puissance de cette dernière par la norme. Plus généralement, on a encore:
\begin{align}
 E^{(i)}\frakg_{i-1}&\equiv0\pmod{\frakg_i},
& E^{(n)}\cdots E^{(i)}\frakg_{i-1}&\equiv0\pmod{\frakm},\notag\\
 R^{(i)}\frakg_{i-1}&\equiv0\pmod{\frakg_i},
& R^{(n)}\cdots R^{(i)}\frakg_{i-1}&\equiv0\pmod{\frakm}, \tag{4}
\end{align}
(H.-N., théorème VIII), et de plus: si $\frakm$ est divisible par $\frakn$, et si les normes $R_\frakm$ et $R_\frakn$ coïncident, alors les idéaux $\frakm$ et $\frakn$ coïncident aussi (H.-N., théorème IX). En remarquant que, pour un diviseur de $\frakm$, la coïncidence de $R^{(1)},R^{(2)},\ldots$ entraîne celle des idéaux fondamentaux $\frakg_1,\frakg_2,\ldots$, et que l'on a toujours $\frakg_0=\frako$, on obtient de manière tout à fait correspondante: si $\frakn$ est un diviseur propre de $\frakm$, alors la première norme individuelle de $\frakn$ qui diffère de celle correspondante de $\frakm$ devient un diviseur propre.\footnote{En revanche, des facteurs ultérieurs de $R_\frakn$ peuvent devenir des multiples des facteurs correspondants de $R_\frakm$; par exemple toujours lorsque $\frakm$ est un idéal premier et que $\frakn$ n'est pas l'idéal unité.} Si $\frakm$ contient un polynôme $G^{(i)}\ne0$, alors, par définition, $\frakg_0,\frakg_1,\ldots,\frakg_{i-1}$ deviennent égaux à l'idéal unité; et ainsi $R^{(1)},\ldots,R^{(i-1)}$, d'après la propriété des normes individuelles indiquée en 5., qui est de déterminer la dimension des quotients correspondants, deviennent égales à l'unité; par conséquent $E^{(1)},\ldots,E^{(i-1)}$ deviennent également égales à l'unité.

\textbf{7. Décomposition des normes individuelles et des diviseurs élémentaires; composantes des idéaux fondamentaux.} Par polynôme irréductible (allemand historique \emph{Primfunktion}) on entend un polynôme irréductible relativement à $P(u)$; par polynôme primaire (\emph{Primärfunktion}), une puissance d'un polynôme irréductible. Un polynôme primaire propre est un polynôme primaire qui n'est pas lui-même irréductible, c'est-à-dire une puissance d'exposant strictement supérieur à un. Une décomposition d'un polynôme en plus grands facteurs primaires est une décomposition où chaque facteur est un polynôme primaire, mais où le produit de deux facteurs quelconques n'est plus primaire.\footnote{La définition de polynôme irréductible et de polynôme primaire pourrait aussi se formuler par analogie exacte avec 11., en remplaçant seulement le mot idéal par le mot polynôme. Les plus grands facteurs primaires sont l'analogue des plus grandes composantes primaires de la décomposition en 12.} À la décomposition de la norme individuelle $R^{(i)}=N(\frakg_{i-1}\mid\frakg_i)$ en plus grands facteurs primaires $Q_\lambda^{(i)}$ correspond une représentation de $\frakg_i$ comme plus petit commun multiple
\[
 \frakg_i=[\frakr_1,\frakr_2,\ldots,\frakr_s],
\]
de telle sorte que $Q_\lambda^{(i)}=N(\frakg_{i-1}\mid\frakr_\lambda)$; $\frakr_\lambda$ possède alors $\frakg_{i-1}$ comme idéal fondamental de $(i-1)$-ième étage, et coïncide avec son idéal fondamental de $i$-ième étage; les $\frakr_\lambda$ sont appelés les composantes des idéaux fondamentaux. Le plus haut diviseur élémentaire de $\frakr_\lambda$ devient égal au plus grand diviseur commun -- au sens polynomial -- de $Q_\lambda^{(i)}$ et de $E^{(i)}$, donc égal à un plus grand facteur primaire de $E^{(i)}$. Les $\frakr_\lambda$ sont déterminés de manière unique par leur comportement indiqué ci-dessus relativement aux idéaux fondamentaux et par l'exigence que la norme individuelle ou le plus haut diviseur élémentaire devienne un plus grand facteur primaire de $R^{(i)}$, respectivement de $E^{(i)}$ (H.-N., théorème X).

\textbf{8. Notion de dimension.} On attribue à l'idéal $\frakm$ la dimension maximale $n-i$ lorsque $R^{(i)}$ est la première norme individuelle différente de l'unité -- ou, ce qui est équivalent d'après 6., lorsque $\frakg_i$ est le premier idéal fondamental différent de l'idéal unité; à l'idéal unité $\frako$ on attribue la dimension $-1$. S'il n'existe qu'une seule norme individuelle $R^{(i)}$ différente de l'unité, donc si $\frakg_i=\frakg_{i+1}=\cdots=\frakm$, alors la dimension maximale est aussi appelée dimension de $\frakm$. Si $\frakm$ contient un polynôme $G^{(i)}\ne0$, il est tout au plus de dimension maximale $n-i$, comme le montre la remarque à la fin de 6. Si $\frakm$ est de dimension maximale $n-i$, et si $\frakn$ est un diviseur de $\frakm$, alors $\frakn$ est donc lui aussi tout au plus de dimension maximale $n-i$.

\textbf{9. Norme (forme résultante) et zéros.} Par zéros d'un idéal $\frakm$ on entend tous les systèmes de valeurs de $x_1,\ldots,x_i$ appartenant à un corps d'extension algébrique convenable de $P(u;x_{i+1},\ldots,x_n)$ ($i=1,2,\ldots,n$), pour lesquels -- après adjonction de $x_{i+1},\ldots,x_n$ au domaine des coefficients -- tous les polynômes de $\frakm$ s'annulent. À chaque facteur $R^{(i)}$ de la norme (forme résultante) correspondent, lors de cette adjonction, autant de zéros de $\frakm$ que l'indique le degré de $R^{(i)}$; et $R^{(i)}$, écrit en combinaisons bilinéaires $w_{\mu\nu}$ des $u_{\mu\nu}$ et d'autres indéterminées $v$, admet la décomposition explicite:
\begin{equation}
 R^{(i)}=\prod\{z-(v_1\bar{x}_{1\nu}+\cdots+v_{i-1}\bar{x}_{i-1,\nu}+v_i\bar{x}_{i,\nu})\}^{\lambda_\nu}, \tag{5}
\end{equation}
où $\bar{x}_{1\nu},\ldots,\bar{x}_{i,\nu}$ désignent un système de zéros associés (H.-N., théorème XIII; une nouvelle fondation en sera donnée au \S{} 3). Parmi les zéros d'un idéal de dimension maximale $n-i$, il y en a donc qui dépendent de $(n-i)$ paramètres, mais aucun qui dépende de plus de paramètres; ainsi est démontrée la concordance de la notion de dimension donnée en 8. avec la formulation usuelle.

\textbf{10. L'anneau de la théorie générale des idéaux.} Le domaine sous-jacent sera un anneau commutatif dans lequel vaut le théorème de la chaîne finie: toute chaîne d'idéaux $\fraka_1,\fraka_2,\ldots,\fraka_k,\ldots$, où chaque $\fraka_i$ est un diviseur propre -- diviseur au sens des idéaux -- du précédent immédiat, s'interrompt au bout d'un nombre fini d'étapes. Cette exigence est équivalente à l'autre, selon laquelle chaque idéal possède une base d'idéal (\emph{Idealtheorie}, théorème I); elle est donc remplie en particulier pour l'anneau de polynômes. Dans les numéros suivants 11., 12., 13., on prendra pour base cet anneau général.

\textbf{11. Idéaux premiers et idéaux primaires.} Un idéal $\frakp$ s'appelle idéal premier si la divisibilité d'un produit par $\frakp$ implique celle d'au moins un facteur; $\frakq$ s'appelle idéal primaire -- ou primaire -- si la divisibilité d'un produit par $\frakq$ implique celle d'un facteur ou d'une puissance de chaque facteur. En symboles:
\[
 a\not\equiv0\pmod{\frakp},\quad b\not\equiv0\pmod{\frakp}\quad\hbox{entraîne}\quad ab\not\equiv0\pmod{\frakp};
\]
\[
 a\not\equiv0\pmod{\frakq},\quad b^x\not\equiv0\pmod{\frakq}\ \hbox{pour toute puissance }x\quad\hbox{entraîne}\quad ab\not\equiv0\pmod{\frakq}.
\]
À tout idéal primaire $\frakq$ correspond un et un seul idéal premier associé $\frakp$, qui est diviseur de $\frakq$ et dont une puissance est divisible par $\frakq$, $\frakq\equiv0\pmod{\frakp}$, $\frakp^\rho\equiv0\pmod{\frakq}$; alors $\frakp$ se compose de tous les éléments de l'anneau dont une puissance est divisible par $\frakq$. Comme le montre la définition, tout idéal premier est en même temps un idéal primaire; par idéaux primaires propres on entend ceux qui ne sont pas en même temps des idéaux premiers, pour lesquels donc $\rho>1$.

\textbf{12. Le théorème de décomposition.} Une représentation $\fraka=[\frakq_1,\ldots,\frakq_s]$ d'un idéal comme plus petit commun multiple s'appelle la plus courte si aucun $\frakq$ n'est contenu dans le plus petit commun multiple des autres -- dans son complément; elle s'appelle représentation par plus grandes composantes primaires si chaque $\frakq$ est primaire, mais si le plus petit commun multiple de deux $\frakq$ quelconques n'est pas primaire. Chaque idéal admet une représentation la plus courte par un nombre fini de plus grandes composantes primaires; pour deux telles représentations différentes, le nombre des composantes et les idéaux premiers associés, tous distincts les uns des autres, coïncident. Les idéaux premiers ainsi déterminés de manière unique seront appelés les idéaux premiers, ou idéaux premiers associés, de l'idéal (\emph{Idealtheorie}, théorème IX).

\textbf{13. Composantes isolées de la décomposition.} Un idéal $\frakr=[\frakq_{i_1},\ldots,\frakq_{i_\lambda}]$ est appelé composante isolée de la décomposition de $\fraka$ lorsque les $\frakq_{i_j}$ interviennent dans au moins une représentation la plus courte de $\fraka$ par plus grandes composantes primaires et satisfont à la condition qu'aucun de leurs idéaux premiers associés ne soit contenu dans l'un des autres idéaux premiers de $\fraka$. Les composantes isolées de la décomposition sont déterminées de manière unique par leurs idéaux premiers; en particulier, les plus grandes composantes primaires isolées sont donc déterminées de manière unique (\emph{Idealtheorie}, théorème XIII).

\textbf{14. Caractéristique, corps parfaits et imparfaits, extension de première espèce.} Un corps $\Omega$ est dit de caractéristique zéro lorsque le corps premier qu'il contient, dérivé de l'unité, est du type du corps des nombres rationnels; il est dit de caractéristique $p$ lorsque ce corps premier est isomorphe à $\mathbb F_p$. Un corps $\Omega$ s'appelle parfait si tout polynôme irréductible à coefficients dans $\Omega$ se décompose en facteurs linéaires distincts dans un corps d'extension convenable; dans le cas contraire, il s'appelle imparfait. Tout corps de caractéristique zéro est parfait; un corps de caractéristique $p$ ne l'est que si et seulement si, avec chaque élément, il contient aussi sa racine $p$-ième. Un polynôme irréductible sur un corps imparfait est un polynôme de degré $r$ en $x^{p^f}$, et se décompose dans un corps d'extension convenable en $p^f$-ièmes puissances de $r$ facteurs linéaires distincts; $f$ est appelé exposant. Un polynôme irréductible est dit de première espèce lorsqu'il se décompose, dans un corps d'extension convenable, en facteurs linéaires distincts, l'exposant $f$ devenant nul; une extension est dite de première espèce lorsque chaque élément devient de première espèce, c'est-à-dire zéro d'un polynôme irréductible de première espèce. Toute extension qui naît par adjonction d'un élément de première espèce est de première espèce; pour les corps parfaits, il n'y a que des extensions de première espèce (Steinitz, \S{} 11 et \S{} 13).

\textbf{15. Degré réduit et exposant d'une extension finie.} Si $\mathcal K$ est une extension finie d'un corps imparfait $\Omega$, alors $\mathcal K$ contient un sous-corps $\mathcal K_0$ de degré $r$, qui devient de première espèce relativement à $\Omega$ et qui se compose de tous et seulement les éléments de première espèce de $\mathcal K$. La $p^f$-ième puissance de tout élément de $\mathcal K$ appartient à $\mathcal K_0$, et il existe dans $\mathcal K$ des éléments qui sont exactement de degré $rp^f$. $r$ est appelé degré réduit, $f$ exposant de l'extension finie (Steinitz, \S{} 14, 1 et \S{} 14, 5).
